\documentclass[aspectratio=169]{beamer}
\usetheme{Madrid}
\usepackage[utf8]{inputenc}
\usepackage[T1]{fontenc}
\usepackage{tikz}
\usepackage{amssymb}
\usepackage{graphicx}
\setbeamertemplate{navigation symbols}{}
\setbeamercolor{title}{fg=white,bg=green!55!teal}
\setbeamercolor{frametitle}{fg=white,bg=green!55!teal}
\setbeamercolor{block title}{bg=green!55!teal!20,fg=black}
\setbeamercolor{block body}{bg=green!55!teal!8,fg=black}
\setbeamercolor{item}{fg=green!55!teal}
\setbeamercolor{itemize item}{fg=green!55!teal}
\setbeamertemplate{itemize item}[triangle]

\title{M5 — Strings \& Validation}
\subtitle{Slicing, methods, checking input — the Caesar cipher in real code}
\author{ByteForge --- Grade 8 --- App Builder}
\date{}

\begin{document}

\begin{frame}
    \titlepage
\end{frame}

\begin{frame}{Today's Plan}
    \begin{itemize}
        \item String Methods \& Slicing
        \item Never Trust Input Blindly
        \item The Caesar Cipher
        \item Reversing a String
        \item .split() on Whitespace
        \item[$\bigstar$] BOSS --- M5 Capstone
    \end{itemize}
\end{frame}

\begin{frame}{.strip() and .lower() — cleaning text}
    \large
    \begin{itemize}
        \item Real user input is messy
        \item \texttt{.strip()} removes leading/trailing spaces, \texttt{.lower()} makes everything lowercase, chaining them does both in one line.
    \end{itemize}
\end{frame}


\begin{frame}{Never trust input blindly — validation}
    \large
    \begin{itemize}
        \item A strong password needs length AND content checks
        \item Combining both with \texttt{and} means BOTH have to be true for acceptance.
    \end{itemize}
\end{frame}


\begin{frame}{The Caesar cipher, in real code}
    \large
    \begin{itemize}
        \item \texttt{ord(ch)} converts a letter to its ASCII number; \texttt{chr(n)} converts back
        \item The exact "text is secretly numbers" idea from Grade 7's ASCII challenge, now doing real work. \texttt{\% 26} wraps the alphabet around.
        \item Worked example: shifting "z" forward by 1 should wrap to "a", not go past the alphabet — \texttt{(25 + 1) \% 26 = 0}, which IS the index for "a". That single \texttt{\% 26} is the whole wraparound trick.
    \end{itemize}
\end{frame}


\begin{frame}{Reversing a string with [::-1]}
    \large
    \begin{itemize}
        \item A slice with no start, no end, and a step of -1 means "walk the whole string, but backward"
        \item Reverses an entire string in one expression.
    \end{itemize}
\end{frame}

\begin{frame}{Real-World Examples}
    \large
    \begin{itemize}
        \item Minecraft's enchanting table language is a real, deliberate obfuscation
        \item Every "Standard Galactic Alphabet" glyph maps to a real English letter, exactly like a Caesar cipher's letter-to-letter mapping, just visual instead of numeric.
    \end{itemize}
\end{frame}


\begin{frame}{Let Them Discover}
    \Large
    Try shifting your own name by 3 letters by hand (A becomes D, Z wraps to C). Where does the wraparound actually happen for you — and does \texttt{\% 26} genuinely handle it the same way?
\end{frame}

\begin{frame}{Your Turn}
    \Large
    Open the student map, find \textbf{M5 — Strings \& Validation}, and work through it
    hands-on --- everything you just saw, they now do themselves.
\end{frame}

\end{document}
